\chapter{The Mathematical Core}

\epigraph{"The thing I enjoy most about math is seeing surprising connections."}{John D. Cook}

In the beginning, the blog was a mixture. Math, code, business, statistics.

By 2018, the mixture had separated. The volatile elements---business advice, quick tips, tech news---had evaporated. What remained was the solute.

The blog became a mathematics blog.

\bigskip

\section*{The Dominance of Math}

The shift is visible in the category data.

\begin{longtable}{p{3.00cm}p{3.00cm}p{3.00cm}}
\toprule
\textbf{Era} & \textbf{Math Posts} & \textbf{\% of Total} \tabularnewline
\midrule
2008-2012 & 357 & 21\% \tabularnewline
2013-2017 & 410 & 42\% \tabularnewline
2018-2025 & 1,611 & 64\% \tabularnewline
\bottomrule
\end{longtable}

In the early years, math was one topic among many. In the later years, it became the atmosphere. By 2020, nearly 70\% of all posts were categorized as Math.

This wasn't a narrowing of interest; it was a deepening. Cook didn't stop writing about computing; he started writing about computing \textit{through} math. He didn't stop writing about statistics; he treated statistics as a branch of applied mathematics.

\bigskip

\section*{The Obsessions}

If you strip away the dates and just look at the topics, certain shapes emerge. These are the obsessions---the subjects Cook returns to, year after year.

\textbf{Pi (over 400 posts).} The most famous constant is also the most frequent guest. Cook writes about calculating pi, finding patterns in pi, the history of pi, and the unexpected places pi appears (like the Mandelbrot set).

\textbf{Primes (103 posts).} The atoms of number theory. The blog tracks the distribution of primes, the gaps between primes, and the algorithms used to find them.

\textbf{Special Functions (140 posts).} This is the blog's signature. The Gamma function, the Bessel functions, the Riemann Zeta function. These are the workhorses of mathematical physics, often neglected in standard curriculums. Cook treats them as old friends.

\begin{blogquote}
"Bessel functions are like the sine and cosine functions, but for cylindrical geometries. They show up whenever you vibrate a drumhead."
\end{blogquote}

\textbf{Fourier Analysis (45 posts).} The mathematics of signals. Cook explores the Fast Fourier Transform (FFT) not just as a tool, but as a mathematical object with beautiful symmetries.

\bigskip

\section*{The Connectors}

The power of the blog isn't just in explaining these topics in isolation. It's in connecting them.

In a single post, Cook might connect:
1. A problem in \textbf{probability} (how likely is a number to be prime?)
2. To a theorem in \textbf{number theory} (the Prime Number Theorem)
3. To a function in \textbf{analysis} (the logarithmic integral)
4. To an implementation in \textbf{Python} (using SciPy)

This is the "surprising connections" he mentions in the epigraph. The blog demonstrates that mathematics is not a collection of separate silos (Algebra, Analysis, Topology) but a single, interconnected web.

A post about "random walks" (probability) leads to "diffusion equations" (differential equations) which leads to "heat kernels" (geometry). The blog walks these paths effortlessly.

\bigskip

\section*{The Library of Functions}

Over seventeen years, the blog has essentially built a library of special functions.

Most programmers know \texttt{sin}, \texttt{cos}, \texttt{exp}, and \texttt{log}.
Cook introduces them to \texttt{erf}, \texttt{gamma}, \texttt{beta}, \texttt{zeta}, and \texttt{bessel}.

He explains what they are, why they matter, and---crucially---how to compute them. He bridges the gap between the abstract definition (an integral from zero to infinity) and the concrete implementation (a function call in Python).

This is why the blog is cited in the documentation for Julia and SciPy. It provides the context that the official docs leave out. It explains the \textit{character} of the functions.

\bigskip

\section*{The Infinite Well}

Why did math win?

Why did clinical trials fade and business advice disappear, while math grew to consume the blog?

Because math is inexhaustible.

You can write about business advice for a few years, but eventually, you repeat yourself. You can write about a specific programming language, but eventually, the language becomes obsolete.

But you can write about the Riemann Zeta function forever. It has infinite depth. It connects to everything. It never gets old.

The move to the mathematical core was a move toward sustainability. By anchoring the blog in the most permanent subject, Cook ensured that he would never run out of things to say.

The blog stopped chasing the news and started exploring the eternal.

\bigskip

\textit{The core was established. The library was built. But in 2019, the quiet world of number theory was interrupted by a sudden, urgent need from the outside world.}
